\documentclass[10pt,a4paper]{article}
\usepackage{amsmath,amssymb,amsthm,graphicx}
% Minimal retained preamble for S3 SIGMA article cards taken from
% sources/sigma-2023-020-planar-orthogonal-polynomials-as-type-i-mult/pinned-original-source.tex
% (doi:10.3842/SIGMA.2023.020, SIGMA 19 (2023), CC BY 4.0).
%
% Omitted and why: the publisher \documentclass{sigma} block and the
% frontmatter commands (FirstPageHeading, ShortArticleName, ArticleName,
% Author, Address, Email, ArticleDates, Abstract, Keywords,
% Classification) - the sigma class is unavailable in the toolchain and
% the assembler supplies the article class; the dsfont package is
% dropped because no counted or context span of this card uses
% \mathds (unused-package disclosure per the three-pitfall broadcast,
% item ii); the TikZ settings block (cross style, tikzdeclarepattern)
% is dropped because no span of this card uses those styles.
% \samepage is a sigma-class spacing command the base LaTeX kernel
% does not define; author spans (source lines 123 and 831) contain it,
% so an empty stub is provided.
% The theorem declarations, equation numbering and paper macros below
% are copied verbatim from the pinned source. Note the counter name is
% Theorem (\newtheorem{Theorem}{Theorem}[section]); initial_counters
% keys use exactly these counter names.
\usepackage{amsmath,amssymb,amsthm}
\usepackage{mathtools}
\usepackage[mathscr]{euscript}
\providecommand{\qed}{\ensuremath{\square}}
\providecommand{\samepage}{}
\allowdisplaybreaks
\numberwithin{equation}{section}
\newtheorem{Theorem}{Theorem}[section]
\newtheorem*{Theorem*}{Theorem}
\newtheorem{Corollary}[Theorem]{Corollary}
\newtheorem{Lemma}[Theorem]{Lemma}
\newtheorem{Proposition}[Theorem]{Proposition}
\newtheorem*{rhproblem}{Riemann--Hilbert problem}
{ \theoremstyle{definition}
 \newtheorem{Definition}[Theorem]{Definition}
 \newtheorem{Note}[Theorem]{Note}
 \newtheorem{Example}[Theorem]{Example}
 \newtheorem{Remark}[Theorem]{Remark}
}
 \DeclareMathOperator{\Leb}{Leb}
 \DeclareMathOperator{\Span}{Span}
 \newcommand*{\Cset}{\mathbb{C}} % complex numbers
 \newcommand*{\Rset}{\mathbb{R}} % real numbers
 \newcommand*{\Nset}{\mathbb{N}} % natural numbers
\title{Equivalent characterisations of planar orthogonal polynomials (Theorem 2.1)}
\author{Berezin, Sergey, Kuijlaars, Arno B.J., Parra, Iván}
\date{doi:10.3842/SIGMA.2023.020; CC BY 4.0}
\begin{document}
\maketitle
\setcounter{section}{2}
\setcounter{Theorem}{0}
\noindent\textit{Editorial note.} All delivered mathematical text is the author's own, copied from the pinned source file sources/sigma-2023-020-planar-orthogonal-polynomials-as-type-i-mult/pinned-original-source.tex; no statement or proof step was rewritten, reordered or removed.\par\medskip
\noindent\textit{Editorial note.} The counted span is the author Theorem environment with its own label; the proof is delivered as two spans - the main proof (source lines 413-545, closing 'completes the proof of Theorem') and the author's alternative proof of (a) implies (c) (source lines 554-604) with its introducing equation (source lines 548-552) as a context span.\par\medskip
\noindent\textit{Editorial note.} The proof takes Q(z) = z\textasciicircum{}k in the fundamental identity (eq:PQscalar); the identity's carrier Proposition prop:Contour1 (printed 3.1) is delivered as a context span placed after the statement, so it prints locally renumbered and the statement keeps 2.1.\par\medskip
\noindent\textit{Editorial note.} Disclosure: the source states the three operator properties W1-W3 (source lines 448-475) without proof, and step (c)-(d) invokes the uniqueness theorem for homogeneous constant-coefficient linear ODEs (source lines 513-514); both are named standard facts the argument reduces to.\par\medskip
\noindent\textit{Editorial note.} Printed number anchored to the published PDF (SIGMA 19 (2023) 020, read from pdftotext of sigma23-020.pdf). The pinned tex has seven sections and no appendix and its section structure matches the published PDF one-to-one, so the unit 010 pinned-versus-published realignment does not trigger.\par\medskip
\noindent\textit{Editorial note.} Equation and theorem numbers inside the card are locally renumbered; every cross-reference is a symbolic \textbackslash{}eqref/\textbackslash{}ref and resolves inside the card.\par\medskip
\noindent\textit{Editorial scope.} Single environment card cut from the SIGMA article doi:10.3842/SIGMA.2023.020 (sigma-2023-020-planar-orthogonal-polynomials-as-type-i-mult); statement and proof are the author's own text at source lines 215-244 and 413-545.\par\medskip
\noindent\textit{Allowed standard prerequisites (editorial).}\begin{itemize}
\item Fundamental identity (Proposition 3.1). For polynomial W with positive integer exponents the W-scalar product equals a contour integral of P(z)W(z) against a ray integral of W*(u)Q*(u) e\textasciicircum{}\{-uz\} (eq:PQscalar, delivered in context).. Load-bearing: step (a)-(b) takes Q(z) = z\textasciicircum{}k in this identity and reads off eq:Pzk.
\item Scalar product and orthogonality system of P\_n. The scalar product of the modified Gaussian measure and the defining vanishing-moment system of P\_n (eq:fgscalar, def:Pn, delivered in context).. Load-bearing: the equivalence is stated for the monic polynomial P\_n of this system; the alternative proof cites def:Pn directly.
\item Operator properties W1, W2, W3. Stated inside the main proof span (source lines 448-475): kernel description eq:Wproperty1, degree behaviour eq:Wproperty2/eq:Wproperty2alt, vanishing-order reduction eq:Wproperty3.. Load-bearing: the (c)-(d) equivalences apply them explicitly; the source states them as basic properties of the product-form operator and proves none of them.
\item Uniqueness theorem for constant-coefficient linear ODEs. A solution of a homogeneous constant-coefficient linear ODE of order c with c vanishing initial conditions is identically zero.. Load-bearing: step (c)-(d) concludes Q = 0 from the c vanishing initial conditions (source lines 513-514); standard named theorem, stated without proof in the source.
\end{itemize}\medskip
 The polynomials in question are orthogonal with respect to a modified Gaussian measure,
 \begin{gather}
 \label{eq:muW}
 \mu_W({\rm d}z) = \frac{1}{\pi} |W(z)|^2 {\rm e}^{-|z|^2} \Leb({\rm d}z),
 \end{gather}
 where $\Leb$ denotes the Lebesgue measure on~$\Cset$ (identified with~$\Rset^2$) and the weight~$W$ reads{\samepage
 \begin{equation}
 \label{eq:Wdef}
 W(z) = \prod_{j=1}^p (z-a_j)^{c_j}, \qquad z \in \Cset,
 \end{equation}
 where~$p \in \Nset$, the~$c_j$ are positive real numbers, and the~$a_j$ are distinct complex numbers.}

 Denote the scalar product corresponding to~\eqref{eq:muW} by
 \begin{equation}
 \label{eq:fgscalar}
 \langle f, g \rangle_W = \int_{\Cset} f(z)\overline{g(z)} \, \mu_W({\rm d}z).
 \end{equation}
 Then, the $n$-th degree monic orthogonal polynomial $P_n$ with
 respect to~$\mu_W$ can be uniquely recovered by solving a linear system of equations for its coefficients,
 \begin{equation}
 \label{def:Pn}
 \int_{\Cset} P_{n}(z)\overline{z}^k\, \mu_W({\rm d}z)=0, \qquad k=0,\dots,n-1.
 \end{equation}

 Our main result is Theorem~\ref{thm:theo1} below. It gives a number of properties that are equivalent to the planar orthogonality corresponding to~\eqref{eq:muW}--\eqref{def:Pn} in the case the~$c_j$ are positive integers. Below, we use the conjugate~$W^*$ of~$W$ defined by
 \begin{equation}
 \label{eq:Wconj}
 W^*(z) = \overline{W(\bar{z})} = \prod_{j=1}^p (z-\overline{a_j})^{c_j}, \qquad z \in \Cset.
 \end{equation}
 We also use~$\mathscr{D}_z$ to denote the derivative operator with respect to~$z$, and then~$W^*(\mathscr{D}_z)$ is the differential operator
 \begin{equation}
 \label{eq:Woperator}
 W^*(\mathscr{D}_z) = \prod_{j=1}^p (\mathscr{D}_z - \overline{a_j})^{c_j}.
 \end{equation}

 \begin{Theorem}
 \label{thm:theo1}
 Let~$W$ be given by~\eqref{eq:Wdef} where all~$c_j$ are positive integers $\big($so that~$W$ is a~polynomial of degree~$c = \sum_{j=1}^p c_j\big)$, and let the~$a_j$ be distinct complex numbers. Then, the following properties are equivalent for a monic polynomial~$P_n$ of degree~$n$,
 \begin{enumerate}
 \item[$(a)$] $P_n$ is the planar orthogonal polynomial on~$\Cset$ with weight~\eqref{eq:muW}.
 \item[$(b)$] $P_n$ satisfies
 \begin{equation}
 \label{eq:Pnequiv2}
 \frac{1}{2\pi {\rm i}} \oint_{\gamma} P_n(z) W(z) \phi_k(z)\, {\rm d}z = 0, \qquad k=0,1, \dots, n-1,
 \end{equation}
 where $\gamma$ is a closed contour around the origin and
 \begin{equation}
 \label{eq:phikdef}
 \phi_k(z) = \int_0^{\bar{z} \times \infty} W^*(u) u^k {\rm e}^{-uz}\, {\rm d}u,
 \end{equation}
 where the path of integration in~\eqref{eq:phikdef} goes from~$0$ to~$\infty$ along the ray $\arg u = \arg \bar{z}$.
 \item[$(c)$] One has
 \begin{equation}
 \label{eq:Pnequiv3}
 W^*(\mathscr{D}_z) \left[ P_n(z) W(z) \right] = O(z^n)
 \end{equation}
 as~$z \to 0$.
 \item[$(d)$] There exist polynomials~$Q_j$ of~$\deg Q_j \leq c_j-1$, for $j=1, \dots, p$, such that
 \begin{equation}
 \label{eq:Pnequiv4}
 P_n(z) W(z) + \sum_{j=1}^p Q_j(z) {\rm e}^{\overline{a_j} z} = O\big(z^{n + c}\big)
 \end{equation}
 as~$z \to 0$.
 \end{enumerate}
 \end{Theorem}

 \begin{Proposition}
 \label{prop:Contour1}
 Let~$P$ and~$Q$ be polynomials and suppose that the~$c_j$ in~\eqref{eq:Wdef} are positive integers.
 Then, the fundamental identity holds,
 \begin{align}
 \langle P, Q \rangle_W &= \frac{1}{\pi} \int_{\Cset} P(z) \overline{Q(z)} |W(z)|^2 {\rm e}^{-|z|^2} \Leb({\rm d}z)\nonumber \\
 & = \frac{1}{2\pi {\rm i}} \oint_{\gamma} P(z) W(z) \int_{0}^{\bar{z} \times \infty} W^*(u) Q^*(u) {\rm e}^{-uz}\, {\rm d}u\, {\rm d}z, \label{eq:PQscalar}
 \end{align}
 where $\gamma$ is a simple closed contour that goes around the origin once in the counterclockwise direction and the path for the~$u$ integral goes from~$0$ to~$\infty$ along the ray~$\arg u = \arg \bar{z}$.
 \end{Proposition}

 \begin{proof}%[Proof: $\rm{(a)} \iff \rm{(b)}$.]
 $(a) \Leftrightarrow (b)$: Taking~$Q(z) = z^k$ in the identity~\eqref{eq:PQscalar} and recalling the definition~\eqref{eq:phikdef} of~$\phi_k$, we obtain
 \begin{equation}
 \label{eq:Pzk}
 \big\langle P, z^k \big\rangle_W = \frac{1}{2\pi {\rm i}} \oint_\gamma P(z) W(z) \phi_k(z)\, {\rm d}z.
 \end{equation}
 This identity \eqref{eq:Pzk} shows that $(a)$ and $(b)$ in Theorem~\ref{thm:theo1}
 are equivalent.

%[Proof: $\rm (b) \iff (c)$.]
\smallskip\noindent $(b) \Leftrightarrow (c)$:
 From basic properties of the Laplace transform, we get
 \begin{equation*}
 %\label{eq:proofeq11}
 \int_{0}^{\bar z \times \infty} W^*(u) u^k {\rm e}^{-uz}\, {\rm d}u = W^*(-\mathscr{D}_z) (-\mathscr{D}_z)^k \bigg[\frac{1}{z}\bigg] = W^*(-\mathscr{D}_z) \bigg[\frac{k!}{z^{k+1}}\bigg],
 \end{equation*}
 Thus, recalling~\eqref{eq:phikdef}, we can write
 \begin{equation*}
 \frac{1}{2\pi {\rm i}} \oint_{\gamma} P_n(z) W(z) \phi_k(z)\, {\rm d}z = \frac{k!}{2\pi {\rm i}} \oint_{\gamma} P_n(z)W(z) W^*(-\mathscr{D}_z) \bigg[\frac{1}{z^{k+1}}\bigg] {\rm d}z
 \end{equation*}
 for any polynomial $P_n$. Integrating by parts, we find
 \begin{equation*}
 %\label{eq:Taylorzero}
 \frac{1}{2\pi {\rm i}} \oint_{\gamma} P_n(z) W(z) \phi_k(z) {\rm d}z = \frac{k!}{2\pi {\rm i}} \oint_{\gamma} \frac{W^*(\mathscr{D}_z)[P_n(z)W(z)]}{z^{k+1}} \, {\rm d}z,
 \end{equation*}
 as there are no boundary terms on the closed contour~$\gamma$ and~$W^*(\mathscr{D}_z)$ is the adjoint of~$W^*(-\mathscr{D}_z)$. Thus, $P_n$ satisfies~\eqref{eq:Pnequiv2} if and only if
 \begin{equation}
 \label{eq:Taylorzero2}
 \frac{1}{2\pi {\rm i}}
 \oint_{\gamma} \frac{W^*(\mathscr{D}_z)[P_n(z)W(z)]}{z^{k+1}} \, {\rm d}z =0, \qquad k =0, \dots, n-1.
 \end{equation}
 The identities~\eqref{eq:Taylorzero2} mean that the coefficient before~$z^k$ of the polynomial~$W^*(\mathscr{D}_z)[P_n(z)W(z)]$ vanishes for~$k=0, \dots, n-1$, and we conclude that items $(b)$ and $(c)$ are equivalent.

 In the proof that $(c)$ and $(d)$ are equivalent, we are going to use three basic properties of the operator~$W^*(\mathscr{D}_z)$ from~\eqref{eq:Woperator}. These properties are stated separately for the ease of reference.

 \medskip\noindent
 \textit{Property $W_1$.} $W^*(\mathscr{D}_z)$ is a linear differential operator with kernel
 \begin{equation}
 \label{eq:Wproperty1}
 \ker W^*(\mathscr{D}_z) = \Bigg\{ \sum_{j=1}^p Q_j(z) {\rm e}^{\overline{a_j} z} \, \bigg|\, \deg{Q_j} \leq c_j-1 \text{ for } j=1, \dots, p \Bigg\}.
 \end{equation}
 The kernel is a vector space over~$\Cset$ and its dimension is~$c = \sum_{j=1}^p c_j$.

 \medskip\noindent
 \textit{Property $W_2$.}
 If~$P$ is a polynomial then~$W^*(\mathscr{D}_z) [P(z)]$ is a polynomial and
 \begin{equation}
 \label{eq:Wproperty2}
 \deg W^*(\mathscr{D}_z) [P(z)] = \deg P(z),
 \end{equation}
 provided all~$a_j \neq 0$. If~$a_1 =0$ and~$\deg P(z) \geq c_1$, then
 \begin{equation}
 \label{eq:Wproperty2alt}
 \deg W^*(\mathscr{D}_z) [P(z)] = \deg P(z)-c_1.
 \end{equation}

 \medskip\noindent
 \textit{Property $W_3$.} $W^*(\mathscr{D}_z)$ reduces the order of vanishing at~$z=0$ by~$c = \deg W$, provided the order of vanishing is greater than or equal to~$c$. Namely, for every non-negative integer~$n$, and analytic function $F$ with $F(z) = O(z^{n+c})$ as $z \to 0$, we have
 \begin{equation}
 \label{eq:Wproperty3}
 W^*(\mathscr{D}_z) [ F(z) ] = O(z^n)
 \end{equation}
 as~$z \to 0$.

 %\begin{proof}[Proof: $\rm (d) \implies (c)$.]
 \smallskip\noindent
 $(d) \Rightarrow (c)$:
 Suppose~$P_n$ satisfies~\eqref{eq:Pnequiv4}. Then, by applying~$W^*(\mathscr{D}_z)$ and using Properties~$W_1$ and~$W_3$, we obtain
 \begin{equation}
 W^*(\mathscr{D}_z) [P_n(z) W(z)] = W^*(\mathscr{D}_z) \big[ O\big(z^{n+c}\big)\big] = O(z^n),
 \end{equation}
 which is~\eqref{eq:Pnequiv3}. Thus, $(d)$ implies $(c)$.

 %[Proof: $\rm (c) \implies (d)$.]
 \smallskip\noindent
$(c) \Rightarrow (d)$:
 In the proof, we assume without loss of generality that~$a_1 =0$ and~$c_1 \geq 0$. This includes the case that~$W$ has no zero at the origin since one can always set~$c_1=0$.

 Consider the linear mapping~$\pi\colon \ker(W^*(\mathscr{D}_z)) \to \Cset^c$ given by
 \begin{equation}
 \label{eq:pidefalt}
 \ker(W^*(\mathscr{D}_z)) \ni Q(z) = \sum_{j=0}^{\infty} q_j z^j \stackrel{\pi}{\longmapsto} (q_0, \dots, q_{c_1-1}, q_{n+c_1}, \dots, q_{n+c-1}) \in \Cset^c.
 \end{equation}
 We claim that~$\pi$ is injective, which will imply that~$\pi$ is surjective since it is a linear mapping between vector spaces of the same dimension~$c$, as follows from~\eqref{eq:pidefalt} and Property~$W_1$. Suppose~$Q \in \ker(W^*(\mathscr{D}_z))$ with~$\pi(Q) = (0, \dots, 0) \in \mathbb C^c$. Then,
 \begin{equation}
 \label{eq:QzOncalt}	
 Q(z) = \sum_{j=c_1}^{n+c_1-1} q_j z^j + O\big(z^{n+c}\big)
 \end{equation}
 as~$z \to 0$.
 From~\eqref{eq:Wproperty3} and~$W^*(\mathscr{D}_z)[Q(z)] = 0$, it then follows that
 \begin{equation}
 \label{eq:WstarOznalt}
 W^*(\mathscr{D}_z)\Bigg[ \sum_{j=c_1}^{n+c_1-1} q_j z^j\Bigg] = O(z^n)
 \end{equation}
 as $z \to 0$. Because of~\eqref{eq:Wproperty2}--\eqref{eq:Wproperty2alt}, we have
 that the left-hand side of~\eqref{eq:WstarOznalt} is a polynomial of degree~$\leq n-1$,
 and due to~\eqref{eq:WstarOznalt} it has a zero at~$z=0$ of order at least~$n$. Hence, \eqref{eq:WstarOznalt} vanishes identically, and thus~$\sum_{j=c_1}^{n+c_1-1} q_j z^j$
 belongs to the kernel of~$W^*(\mathscr{D}_z)$. By Property~$W_1$, the kernel contains polynomials up to degree $c_1-1$ but no polynomials of higher degrees. Hence, $q_j=0$ for $j=c_1, \dots, n+c_1-1$. Consequently, by~\eqref{eq:QzOncalt}
 we have that~$Q(z) = O(z^{n+c})$ as~$z \to 0$.
 In particular, $Q(0) = Q'(0) = \cdots = Q^{(c-1)}(0) = 0$.
 Thus~$Q \in \ker(W^*(\mathscr{D}_z))$ is a solution of a~homogeneous constant coefficient linear ODE of order $c$ with $c$ vanishing initial conditions.
The uniqueness theorem for such ODEs then yields~$Q \equiv 0$, which justifies the claim.

 Now, assume that~$P_n$ satisfies~\eqref{eq:Pnequiv3}, and write
 \begin{equation*}
 P_n(z) W(z) = \sum_{j=0}^{n+c} p_j z^j.
 \end{equation*}
 Since~$\pi$ is surjective, there is~$Q \in \ker(W^*(\mathscr{D}_z))$,
 \begin{equation*}
 Q(z) = \sum_{j=0}^{\infty} q_j z^j,
 \end{equation*}
 such that
$
 \pi(Q) = (-p_0, \dots, -p_{c_1-1}, -p_{n+c_1}, \dots, -p_{n+c-1}).
$
 Then,
 \begin{equation}
 \label{eq:PnWQalt}
 P_n(z) W(z) + Q(z) = \sum_{j=c_1}^{n+c_1-1} (p_j + q_j) z^j + O\big(z^{n+c}\big).
 \end{equation}
 Applying~$W^*(\mathscr{D}_z)$ to~\eqref{eq:PnWQalt} and using~\eqref{eq:Pnequiv3}, $W^*(\mathscr{D}_z)[Q(z)]=0$, and~\eqref{eq:Wproperty3}, we obtain
 \begin{equation}
 \label{eq:Wpjqjalt}
 W^*(\mathscr{D}_z)\Bigg[\sum_{j=c_1}^{n+c_1-1} (p_j + q_j) z^j\bigg] = O(z^{n})
 \end{equation}
 as $z \to 0$. We find ourselves in the situation similar to that while proving the injectivity of~$\pi$. It followed from~\eqref{eq:WstarOznalt} that~$q_j = 0$ for~$j=c_1,\dots, n+c_1-1$. In the exactly same way, now it follows from~\eqref{eq:Wpjqjalt} that~$p_j+q_j = 0$ for~$j=c_1, \dots, n+c_1-1$.

 Due to~\eqref{eq:PnWQalt}, we obtain
 \begin{equation*}
 P_n(z) W(z) + Q(z) = O\big(z^{n+c}\big),
 \end{equation*}
 which is exactly~\eqref{eq:Pnequiv4} because~$Q \in \ker(W^*(\mathscr{D}_z))$ must be of the form given in~\eqref{eq:Wproperty1}. This shows that $(c)$ implies $(d)$ and completes the proof of Theorem~\ref{thm:theo1}.
 \end{proof}

 We present an alternative proof of~$(a) \Rightarrow (c)$ by making use of the Fourier transform of the measure $\mu_W$ in~\eqref{eq:muW},
 \begin{equation}
 \label{eq:muhat}
 \widehat{\mu}_W(\zeta, \overline{\zeta}) = \int_{\Cset} {\rm e}^{\overline{\zeta} z + \zeta \bar{z}} \mu_W({\rm d}z).
 \end{equation}

 \begin{proof}[Alternative proof]
 $(a) \Rightarrow (c)$:
 Apply the Wirtinger derivatives
 \begin{equation*}
 \mathscr{D}_{\zeta} = \frac{1}{2}\bigg(\frac{\partial}{\partial u} - {\rm i} \frac{\partial}{\partial v}\bigg), \qquad
 \mathscr{D}_{\overline{\zeta}} = \frac{1}{2}\bigg(\frac{\partial}{\partial u} + {\rm i} \frac{\partial}{\partial v}\bigg),
 \end{equation*}
 where~$\zeta = u + {\rm i}v$, to~\eqref{eq:muhat}, and observe that
 \begin{equation}
 \label{eq:Dmuhat}
 \mathscr{D}_{\overline{\zeta}}^j \mathscr{D}_{\zeta}^k \big[\widehat{\mu}_W(\zeta, \overline{\zeta})\big] = \int_{\Cset} z^j \overline{z}^k {\rm e}^{\overline{\zeta} z + \zeta \overline{z}} \mu_W({\rm d}z).
 \end{equation}

 Due to linearity, \eqref{eq:Dmuhat} leads to
 \begin{equation*}
 P_n(\mathscr{D}_{\overline{\zeta}}) \mathscr{D}_{\zeta}^k \big[ \widehat{\mu}_W(\zeta, \overline{\zeta}) \big] = \int_{\Cset} P_n(z) \overline{z}^k {\rm e}^{\overline{\zeta} z + \zeta \overline{z}} \mu_W({\rm d}z).
 \end{equation*}
 Assuming that~$P_n$ is the planar orthogonal polynomial satisfiying~\eqref{def:Pn}, we have
 \begin{equation}
 \label{eq:remeq2}
 P_n(\mathscr{D}_{\overline{\zeta}}) \mathscr{D}_{\zeta}^k \big[ \widehat{\mu}_W(\zeta,\overline{\zeta}) \big] \big|_{\zeta=0, \overline{\zeta} = 0}= 0, \qquad k=0,\dots,n-1,
 \end{equation}
 where we view $\zeta$ and $\overline{\zeta}$ as two independent variables. The latter is permissible as long as the Fourier transform~$\widehat{\mu}_W(\cdot,\cdot)$, as a function of two complex variables, extends analytically to the neighborhood of~$\big\{(\xi_1,\xi_2) \in \Cset \times \Cset\mid \xi_2=\overline{\xi_1}\big\}$. This is the case we encounter below.

 For~$\mu_W$ from~\eqref{eq:muW} with polynomial~$W$, the Fourier transform can be computed as follows,
 \begin{align}
 \label{eq:hatmuW}
 \widehat{\mu}_W(\zeta,\bar{\zeta}) &=\frac{1}{\pi} \int_{\Cset} {\rm e}^{ \overline{\zeta} z + \zeta \overline{z}} |W(z)|^2 {\rm e}^{-|z|^2} \Leb{({\rm d}z)}\nonumber
 \\
 &= \frac{1}{\pi} W(\mathscr{D}_{\overline{\zeta}})W^*(\mathscr{D}_\zeta) \int_{\Rset^2} {\rm e}^{2(u x + v y)} {\rm e}^{-x^2 -y^2}\, {\rm d}x\,{\rm d}y \qquad
\nonumber
 \\
 &= W(\mathscr{D}_{\overline{\zeta}})W^*(\mathscr{D}_\zeta) \big[{\rm e}^{\zeta \overline{\zeta}} \big],
 \end{align}
 where $\zeta = u + {\rm i}v$, $z=x+ {\rm i}y$.

 Inserting~\eqref{eq:hatmuW} into~\eqref{eq:remeq2}, and changing the order of the differential operators, we obtain
 \begin{equation}
 \label{eq:remeq3}
 \mathscr{D}_{\zeta}^k W^*(\mathscr{D}_\zeta) (P_n W)(\mathscr{D}_{\overline{\zeta}}) \big[{\rm e}^{\zeta \overline{\zeta}}\big] \big|_{\zeta=0, \bar{\zeta} =0}= 0, \qquad k=0,\dots,n-1.
 \end{equation}
 Since
 \begin{equation*}
 (P_n W)(\mathscr{D}_{\overline{\zeta}}) \big[{\rm e}^{\zeta \overline{\zeta}} \big]	\big|_{\overline{\zeta}=0} = P_n(\zeta) W(\zeta),
 \end{equation*}
 we see that~\eqref{eq:remeq3} implies
 \begin{equation*}
 \mathscr{D}_{\zeta}^k W^*(\mathscr{D}_\zeta) [P_n(\zeta)W(\zeta)] \big|_{\zeta=0} = 0,\qquad k=0,\dots,n-1,
 \end{equation*}
 which is equivalent to~\eqref{eq:Pnequiv3}.
 \end{proof}

\end{document}
